% In this file you should put all LaTeX macros and settings to be used both by
% the pdf version and the web version.
% This should be most of your macros.

% The theorem-like environments defined below are those that appear by default
% in the dependency graph. See the README of leanblueprint if you need help to
% customize this.
% The configuration below use the theorem counter for all those environments
% (this is what the [theorem] arguments mean) and never resets it.
% If you want for instance to number them within chapters then you can add
% [chapter] at the end of the next line.
 % THEOREMS

\newtheorem{theorem}{Theorem}
\newtheorem{remark}{Remark}
\newtheorem{proposition}[theorem]{Proposition}
\newtheorem{observation}{Observation}
\newtheorem{lemma}[theorem]{Lemma}
\newtheorem{corollary}[theorem]{Corollary}
\newtheorem{example}{Example}
\newtheorem{fact}{Fact}
\newtheorem{conjecture}{Conjecture}

\theoremstyle{definition}
\newtheorem{definition}[theorem]{Definition}

% QI MACROS

\newcommand{\braket}[2]{\mbox{$ \langle #1 | #2 \rangle $}}
\newcommand{\sandwich}[3]{\mbox{$ \langle #1 | #2 | #3 \rangle $}}
\newcommand{\ket}[1]{\mbox{$ | #1 \rangle $}}
\newcommand{\bra}[1]{\mbox{$ \langle #1 | $}}
\newcommand{\ev}[1]{\langle #1 \rangle}
\newcommand{\dyad}[1]{\rvert #1 \rangle \! \langle #1 \lvert}
\newcommand{\ketbra}[2]{\lvert #1 \rangle \! \langle #2 \rvert}
\newcommand{\expval}[2]{\langle #2 \lvert #1 \rvert  #2 \rangle}
\newcommand{\matrixel}[3]{\langle #1 \lvert #2 \rvert  #3 \rangle}
\newcommand{\commu}[2]{\left[ #1,\, #2 \right]}
\newcommand{\kett}[1]{\mbox{$ | #1 \rangle \rangle$}}
\newcommand{\braa}[1]{\mbox{$ \langle \langle #1 | $}}

% MATH MACROS
\newcommand{\id}{\mathbb{I}}
\newcommand{\tr}{\mathop{\mathrm{Tr}}}
\newcommand{\mc}{\mathcal}

\newcommand{\norm}[1]{\left\lVert #1 \right\rVert}
\newcommand{\abs}[1]{\left\lvert #1 \right\rvert}

\newcommand{\rr}{\mathbb{R}}
\newcommand{\cc}{\mathbb{C}}

%%%%%%%%%%%%%%%%%%%%%%%%%%%%%%%%%
\newcommand{\equref}[1]{Eq.~\eqref{#1}}
\newcommand{\equsref}[2]{Eqs.~(\ref{#1}) and (\ref{#2})}
\newcommand{\figref}[1]{Fig.~\ref{#1}}
\newcommand{\supfigref}[1]{Supplementary Fig.~\ref{#1}}
\newcommand{\secref}[1]{Sec.~\ref{#1}}
\newcommand{\refcite}[1]{Ref.~\cite{#1}}
\newcommand{\refscite}[1]{Refs.~\onlinecite{#1}}
\newcommand{\refsacite}[2]{Refs.~\onlinecite{#1} and \onlinecite{#2}}
\newcommand{\tableref}[1]{Table~\ref{#1}}
\newcommand{\appref}[1]{Supplementary Materials~\ref{#1}}
\newcommand{\defref}[1]{Definition~\ref{#1}}
\newcommand{\thmref}[1]{Theorem~\ref{#1}}
\newcommand{\lemref}[1]{Lemma~\ref{#1}}
\newcommand{\exref}[1]{Example~\ref{#1}}
\newcommand{\corref}[1]{Corollary~\ref{#1}}
\newcommand{\obsref}[1]{Observable~\ref{#1}}
\newcommand{\propref}[1]{Proposition~\ref{#1}}
